%! Justifying a Claim Based on a Confidence Interval for a Population Proportion — Unit 6, Lesson 6.3 — Homework
\documentclass[10pt]{article}
\usepackage{apstats-article}
\usepackage{apstats-boxes}
\usepackage{grading}
\gradingassignment{apstats-u06-l03-hw}
\gradingcourse{advanced_statistics}
\newcommand{\TallMath}[1]{$\displaystyle #1\rule[-1.4em]{0pt}{3.2em}$}

\begin{document}

\pageheader{Unit 6, Lesson 6.3}{Homework}

\vspace{0.15in}

\begin{enumerate}

  \item A 95\% CI for the proportion of adults who exercise at least 3 days per week is
        $(0.42, 0.54)$.
        \begin{enumerate}[label=\alph*.]
          \item Write a complete interpretation of this CI. \gradedhide{\writeline\\\gradedhere{\gradedblank[id=hw1a, type=rubric, points=2, display={95\% confident (0.42, 0.54) captures the true proportion of all adults who exercise 3+ days a week.}]{\linewidth+6pt}{\gradedlinesheight{1}}}\writelines{9}}
          \item A fitness organization claims 50\% of adults exercise at least 3 days per week.
                Is this plausible? \gradedhide{\writeline\\\gradedhere{\gradedblank[id=hw1b, type=rubric, points=2, display={Yes: 0.50 is inside (0.42, 0.54), so 50\% is plausible.}]{\linewidth+6pt}{\gradedlinesheight{1}}}\writelines{9}}
          \item The surgeon general claims \emph{fewer than} 45\% of adults exercise regularly.
                Does the CI support this claim? \gradedhide{\writeline\\\gradedhere{\gradedblank[id=hw1c, type=rubric, points=2, display={No: the interval includes values above 0.45 (and 0.45 itself), so it does not support $p < 0.45$.}]{\linewidth+6pt}{\gradedlinesheight{1}}}\writelines{9}}
        \end{enumerate}

  \newpage
  \item A 99\% CI for the proportion of students who own a calculator is $(0.71, 0.89)$.
        \begin{enumerate}[label=\alph*.]
          \item A teacher claims 90\% of students own a calculator. Plausible? \gradedhide{\writeline\\\gradedhere{\gradedblank[id=hw2a, type=rubric, points=2, display={No: 0.90 is outside (0.71, 0.89), so the claim is not plausible.}]{\linewidth+6pt}{\gradedlinesheight{1}}}\writelines{9}}
          \item Would a 95\% CI using the same data be wider or narrower? Explain.
                \gradedhide{\writeline\\\gradedhere{\gradedblank[id=hw2b, type=rubric, points=2, display={Narrower: smaller $z^*$ (1.960 vs.\ 2.576), so a smaller margin of error.}]{\linewidth+6pt}{\gradedlinesheight{1}}}\writelines{9}}
          \item If the sample size doubled (holding $\hat p$ constant), how would the width of
                the 99\% CI change? \gradedhide{\writeline\\\gradedhere{\gradedblank[id=hw2c, type=rubric, points=2, display={Narrower: width is multiplied by $1/\sqrt{2} \approx 0.71$, since $ME \propto 1/\sqrt{n}$.}]{\linewidth+6pt}{\gradedlinesheight{1}}}\writelines{9}}
        \end{enumerate}

  \newpage
  \item \textbf{Correcting errors.} For each student statement, identify and correct the error.
        \begin{enumerate}[label=\alph*.]
          \item ``There is a 95\% probability that the true proportion of left-handed students
                is between 0.09 and 0.17.'' \gradedhide{\writeline\\\gradedhere{\gradedblank[id=hw3a, type=rubric, points=2, display={Not a probability: we are 95\% confident (0.09, 0.17) captures the true proportion of left-handed students.}]{\linewidth+6pt}{\gradedlinesheight{1}}}\writelines{9}}
          \item ``95\% of all students are left-handed according to our interval.'' \gradedhide{\writeline\\\gradedhere{\gradedblank[id=hw3b, type=rubric, points=2, display={95\% is the confidence level: we are 95\% confident the true proportion is in (0.09, 0.17).}]{\linewidth+6pt}{\gradedlinesheight{1}}}\writelines{9}}
        \end{enumerate}

\end{enumerate}

\begin{reflectionbox}
\textbf{Preview:} Lesson 6.4 introduces a new inference tool --- the \emph{significance test} for
a population proportion. Instead of estimating $p$ with an interval, we will ask: is the sample
evidence strong enough to reject a specific claimed value of $p$? Think about how this differs
from saying a value is ``not plausible.''
\end{reflectionbox}

\end{document}
